
\chapter{绪论}

本模板参照南开大学学位论文写作规范编写，目标是在尽量少修改类文件的前提下，为用户提供一套可以直接开始写作的 \LaTeX 学位论文模板。

\section{模板定位}

模板当前主要面向中文论文写作，默认使用 XeLaTeX 编译，并配合 biber 处理参考文献。硕士、博士论文为主要适配对象，本科部分也提供了单独入口和基础版式支持，方便在统一仓库中维护不同学位类型的示例。

与单纯展示排版效果不同，这份示例内容更强调“如何上手使用模板”。因此，后续章节会分别介绍：

\begin{itemize}
    \item 如何理解项目结构与编译方式；
    \item 如何替换模板中的个人信息与正文内容；
    \item 如何使用图片、公式、代码、表格、参考文献等常见学术写作元素；
    \item 如何启用打印模式、匿名评审页、本科双学位封面等可选功能。
\end{itemize}

\section{使用建议}

第一次使用时，建议按照“选择入口文件 $\rightarrow$ 填写 \texttt{\textbackslash nkusetup\{\}} $\rightarrow$ 修改摘要与正文 $\rightarrow$ 整理参考文献”的顺序进行。这样可以先确保整体结构编译通过，再逐步替换为自己的内容。

如果你写的是硕博论文，请优先修改 \texttt{main.tex}；如果你写的是本科论文，请直接使用 \texttt{main-bachelor.tex}。正文示例文件统一放在 \texttt{data/} 目录中，便于按章节逐个替换。

\section{项目地址与反馈}

项目地址为：\url{https://github.com/zhasion/nkuthesis}

如果使用过程中发现模板存在格式错误、文档歧义或兼容性问题，欢迎通过仓库 Issue 反馈并补充说明。若只是一般性的 \LaTeX 写作问题，也建议先结合搜索引擎、官方文档和社区资料排查。

\textbf{免责声明：} 本模板旨在为论文写作提供便利，但作者不保证其完全符合学校或各学院的全部细节要求，亦不对因使用本模板所产生的任何后果承担责任。正式提交前，请务必结合所在学院要求自行核对。
